%!TEX root = ../template.tex
%%%%%%%%%%%%%%%%%%%%%%%%%%%%%%%%%%%%%%%%%%%%%%%%%%%%%%%%%%%%%%%%%%%%
%% chapter7.tex
%% NOVA thesis document file
%%
%% Chapter 7: Conclusion
%%%%%%%%%%%%%%%%%%%%%%%%%%%%%%%%%%%%%%%%%%%%%%%%%%%%%%%%%%%%%%%%%%%%

\typeout{NT FILE chapter7.tex}%

\chapter{Conclusion}
\label{cha:conclusion}

This thesis asked how software modelling is experienced, taught, and learned in higher education, treating modelling as a socio-technical activity rather than as notation drill alone. Fourteen interviews (four lecturers, ten students) were analysed with a lightweight STGT procedure: open coding, short memos, constant comparison, and a findings map. The result is an empirically grounded account of \textbf{14 categories} and \textbf{four phenomena}, sufficient for the research questions of this MSc, not a finished global theory.

\section{Answers to the research questions}
\label{sec:conc_rqs}

On RQ1, the main challenges are a gated first modelling move, symbol and framework overload as a distinct problem, difficulty judging ``good enough'' without a compile result, and breakdowns when students switch notation families. The four phenomena in Chapter~\ref{cha:findings} do not replace those categories: they name the cross-cutting stories that several categories feed (Figures~\ref{fig:findings_map} and~\ref{fig:relations_map}).

On RQ2, teacher examples remain the principal unlock, students invent process tactics when teaching is thin, feedback is both late and unequally accessible, and AI---where used---explains and compares rather than generating graded models.

On RQ3, semantic checks are wanted but tool friction is costly, groups and real audiences decide whether diagrams survive, and course load blocks the demonstrations and loops that the other categories require.

\section{Contributions}
\label{sec:conc_contributions}

The contribution is a bounded, traceable map of modelling-education challenges and support needs that connects student and lecturer voices to the socio-technical conditions of courses, tools, groups, and feedback. Relative to method-specific error taxonomies and Bloom-based learning-outcome frameworks, the account splits mechanisms that are often collapsed (start-gate versus symbol overload; feedback timing versus access; checking wanted versus tools dropped) and treats family translation as its own challenge. Implications for teaching and tools follow the phenomena: keep a worked path into the first diagram; supply a correctness signal while the artefact can still change and in a format students will use; prefer checks and explanations over generators; teach bridges between families; keep an owner and a reader in the loop. Those implications were not trialled here.

\section{Limits}
\label{sec:conc_limits}

Sampling stopped when wave~2 densified properties without adding a 15th overview category, and when further interviews were no longer realistic. That is local saturation of the category set plus a recruitment constraint---not global theoretical saturation. Lecturer-side categories were not re-sampled. Several open questions remain labelled for later work (Table~\ref{tab:open_questions}). Transferability is analytical: the conditions under which a category appeared, not a statistical claim about modelling courses in general.

\section{Closing}
\label{sec:conc_closing}

Future work can take up the unfulfilled questions---matched course pairs, syllabus-level translation teaching, non-leading tests of abstraction, and design trials of private checkers---without discarding them or claiming they were closed. The recipes are in Section~\ref{sec:disc_future}; they are a different project from the map this thesis has already earned.
